% !TeX program = xelatex
\documentclass[12pt]{nextbook}

\UseTemplateSet{
  fonts = {libertinus},
  features = {hyperlinks}
}

\usepackage{paracol}
\usepackage{xcolor}
\usepackage{eso-pic}
\usepackage{needspace}

\vbadness=10000
\hbadness=10000
\hfuzz=3pt
\emergencystretch=1.5em

\makeatletter
\AtBeginDocument{%
  \@ifpackageloaded{hyperref}{%
    \pdfstringdefDisableCommands{%
      \def\ruby#1#2{#1}%
      \def\ReaderVocabText#1{#1}%
      \def\ReaderVocabTerm#1{#1}%
    }%
  }{}%
}
\makeatother

\providecommand{\ReaderColumnSep}{6mm}
\providecommand{\ReaderColumnRatio}{0.75}
\providecommand{\ReaderVocabBaselineskip}{20pt}
\providecommand{\ReaderVocabText}[1]{\textbf{#1}}
\providecommand{\ReaderVocabTerm}[1]{\textbf{#1}}
\providecommand{\ReaderVocabSeparator}{:}
\providecommand{\ReaderVocabItem}[2]{\noindent\ReaderVocabTerm{#1}\ReaderVocabSeparator\ #2\par}

\setlength{\columnsep}{\ReaderColumnSep}
\setlength{\columnseprule}{0pt}
\columnratio{\ReaderColumnRatio}
\swapcolumninevenpages

\newcounter{readersection}
\newwrite\vocabfile
\newif\iffirstreaderlabelwritten
\newlength{\readercolumnswidth}
\newlength{\readertitlegutter}
\setlength{\readercolumnswidth}{\dimexpr\textwidth-\columnsep\relax}
\setlength{\readertitlegutter}{.75\columnsep}

\makeatletter
\newcommand{\readerseparatorline}{%
  \if@mainmatter
    \AtTextLowerLeft{%
      \ifodd\value{page}%
        \put(\LenToUnit{\dimexpr .75\readercolumnswidth+.5\columnsep\relax},0){%
          \color{black!35}\rule{0.3pt}{\textheight}%
        }%
      \else
        \put(\LenToUnit{\dimexpr .25\readercolumnswidth+.5\columnsep\relax},0){%
          \color{black!35}\rule{0.3pt}{\textheight}%
        }%
      \fi
    }%
  \fi
}
\makeatother

\AddToShipoutPictureBG{\readerseparatorline}

\newcommand{\currentvocabfile}{\jobname-vocab-\thereadersection.voc}
\newcommand{\beginvocabwrite}{\immediate\openout\vocabfile=\currentvocabfile}
\newcommand{\finishvocabwrite}{\immediate\closeout\vocabfile}
\newcommand{\vocabitem}[2]{\ReaderVocabItem{#1}{#2}}
\newcommand{\printvocablist}{%
  \footnotesize
  \setlength{\parindent}{0pt}%
  \setlength{\parskip}{0pt}%
  \setlength{\baselineskip}{\ReaderVocabBaselineskip}%
  \raggedright
  \input{\currentvocabfile}%
}

\newcommand{\readersectiontitle}[1]{%
  \par
  \begingroup
  \ifodd\value{page}%
    \leftskip=\readertitlegutter%
    \rightskip=0.25\textwidth plus 1fil%
  \else
    \leftskip=\dimexpr .25\readercolumnswidth+.5\columnsep+\readertitlegutter\relax%
    \rightskip=0pt plus 1fil%
  \fi
  \section{#1}%
  \par
  \endgroup
}

\NewDocumentCommand{\voc}{omm}{%
  \ReaderVocabText{#2}%
  \IfNoValueTF{#1}{%
    \immediate\write\vocabfile{\string\vocabitem{#2}{#3}}%
  }{%
    \immediate\write\vocabfile{\string\vocabitem{#1}{#3}}%
  }%
}

\newenvironment{readersection}[1]{%
  \stepcounter{readersection}%
  \beginvocabwrite
  \Needspace{10\baselineskip}%
  \iffirstreaderlabelwritten\else
    \label{reader:firstpage}%
    \global\firstreaderlabelwrittentrue
  \fi
  \readersectiontitle{#1}%
  \begin{paracol}{2}
}{%
  \finishvocabwrite
  \switchcolumn
  \printvocablist
  \end{paracol}
}

\usepackage{polyglossia}
\setmainlanguage{german}
\setotherlanguage{english}

\title{Deutsches Lesebuch für Anfänger}
\author{Kelvin Yang}
\date{\today}

\begin{document}
\frontmatter
\maketitle
\tableofcontents
\mainmatter
\chapter{Grundtexte}
\begin{readersection}{Der Arzt, der Berliner und die Wahrheit}

Ich bin ein \voc[der Berliner]{Berliner}{Berliner}. Ich habe ein \voc[das Kind]{Kind}{child}, eine \voc[die Frau]{Frau}{woman} und einen \voc[gut]{guten}{good} Freund. Der Freund ist \voc[der Arzt]{Arzt}{doctor}. Er ist ein guter \voc[der Mann]{Mann}{man} und ein guter \voc[der Mensch]{Mensch}{person}. Er hat ein \voc[schön]{schönes}{beautiful} \voc[das Gesicht]{Gesicht}{face} und ein gutes \voc[das Herz]{Herz}{heart}. Viele Leute \voc[sagen]{sagen}{to say}: „Der Arzt \voc[helfen]{hilft}{to help} viel.“ Auch die Frau sagt: „Er hilft, und die \voc[die Natur]{Natur}{nature} \voc[heilen]{heilt}{to heal}.“

Das Kind ist oft krank. Es \voc[essen]{ißt}{to eat} nicht viel, und es \voc[sprechen]{spricht}{to speak} nicht viel. Ich \voc[denken]{denke}{to think}: „Das Kind \voc[brauchen]{braucht}{to need} den Arzt.“ Die Frau \voc[bringen]{bringt}{to bring} das Kind zum Arzt. Der Arzt \voc[sehen]{sieht}{to see} das Kind, \voc[nehmen]{nimmt}{to take} die Hand des Kindes und sagt: „Ich helfe.“ Das Kind \voc[geben]{gibt}{to give} keine \voc[die Antwort]{Antwort}{answer}. Es \voc[scheuen]{scheut}{to avoid} den Arzt, aber der Arzt ist gut.

Der Arzt \voc[lesen]{liest}{to read} ein \voc[das Wort]{Wort}{word} in einem Buch. Er sagt: „Das Herz ist gut. Die Natur heilt. Aber der Mensch braucht auch Zeit.“ Die Frau fragt: „Was hat das Kind?“ Der Arzt sagt die \voc[die Wahrheit]{Wahrheit}{truth}: „Das Kind \voc[wachsen]{wächst}{to grow}, aber es ist schwach. Es braucht Ruhe.“ Ich \voc[wissen]{weiß}{to know} nicht alles, aber ich weiß: Der Arzt \voc[lügen]{lügt}{to lie} nicht.

Der \voc[der Richter]{Richter}{judge} kommt auch. Er ist ein Freund des Arztes. Der Richter sagt: „Ein gutes Wort \voc[finden]{findet}{to find} einen guten \voc[der Ort]{Ort}{place}.“ Der \voc[der Narr]{Narr}{fool} hört das \voc[das Gespräch]{Gespräch}{conversation} und lacht. Er sagt: „Der Mensch ist, was er ißt!“ Die Leute lachen auch. Der Arzt sagt: „Narr, du sagst oft gute Worte, aber du \voc[vergessen]{vergißt}{to forget} auch viel.“ Der Narr \voc[waschen]{wäscht}{to wash} sein Gesicht und ißt eine \voc[die Bratwurst]{Bratwurst}{sausage}.

Es \voc[regnen]{regnet}{to rain}. Die Frau \voc[machen]{macht}{to make} \voc[das Feuer]{Feuer}{fire}. Das Kind sieht das Feuer und \voc[werden]{wird}{to become} ruhig. Der Arzt sagt: „Das Feuer hilft nicht, aber es ist schön.“ Die Frau sagt: „Das Gespräch ist gut, aber zu lang.“ Der Arzt \voc[kürzen]{kürzt}{to shorten} das Gespräch und sagt nur: „Das Kind braucht Ruhe.“

Am \voc[der Kirchhof]{Kirchhof}{graveyard} steht ein Mann. Er ist alt und sagt nichts. Ich sehe den Mann und denke: „Der \voc[der Weg]{Weg}{way} ist kurz, und das Leben ist nicht immer schön.“ Der Arzt sagt: „Der Weg ist kurz, aber das Herz kann gut sein.“ Die Frau gibt dem Kind ein Wort, und das Wort ist gut. Das Kind schläft.

Am Abend \voc[lassen]{lassen}{to let} wir das Feuer brennen. Die Frau sagt: „Wir haben Hoffnung.“ Ich sage: „Ja. Der Arzt hilft, die Natur heilt, und das Kind wird stark.“ Der Berliner lacht, der Narr tanzt, der Richter spricht nicht, und die Bratwurst ist weg.

Der Arzt sagt zu den Kindern: „Ihr braucht Ruhe.“ Dann sagt er zur Frau: „Sie haben ein gutes Kind.“ Ein \voc[neu]{neuer}{new} Tag kommt. Das Kind sieht ein \voc[gebrannt]{gebranntes}{burned} Stück Brot und lacht. Der Narr sagt: „Das Brot \voc[brechen]{bricht}{to break}!“ Der Arzt lacht auch.

\end{readersection}

\begin{readersection}{Der Bericht vom Hof}

Am Freitag kommt ein Student aus der Stadt auf den Hof des \voc[der Bauer]{Bauern}{farmer}. Er will einen Bericht über die \voc[die Arbeit]{Arbeit}{work} auf dem Land schreiben. Der Hof liegt hinter einem kleinen \voc[das Haus]{Haus}{house}. Vor dem Haus steht eine \voc[die Kuh]{Kuh}{cow} mit einer \voc[die Ziege]{Ziege}{goat} und einem jungen \voc[das Kalb]{Kalb}{calf}. Neben der Tür liegt ein \voc[der Kieselstein]{Kieselstein}{pebble}; unter dem Tisch im Hof liegt noch ein Stück \voc[das Eis]{Eis}{ice}, denn die Nacht war kalt.

Der Bauer heißt Martin. Er ist ein ruhiger \voc[der Mann]{Mann}{man}, aber er ist auch \voc{fleißig}{industrious}. Seine \voc[die Mutter]{Mutter}{mother} wohnt im Haus, und seine \voc[die Tochter]{Tochter}{daughter} Anna hilft ihm. Anna trägt am Morgen ein blaues \voc[das Kleid]{Kleid}{dress}, aber sie zieht bald alte Schuhe an, denn der Boden ist naß. Der Student sieht ihre \voc[der Fuß]{Füße}{foot} und ihre nassen \voc[das Bein]{Beine}{leg}. Anna lacht und sagt: „Auf einem Hof \voc[bleiben]{bleibt}{to remain} kein Kleid \voc[lang]{lange}{long} schön.“

Der Tag beginnt ohne \voc[die Hast]{Hast}{haste}. Am \voc[der Anfang]{Anfang}{beginning} \voc[zeigen]{zeigt}{to show} der Bauer dem Studenten den Stall, die Kuh, die Ziege und das Kalb. Dann \voc[machen]{macht}{to do, make} Anna eine kleine \voc[die Übung]{Übung}{practice}: Sie gibt dem Kalb Wasser, sie \voc[bringen]{bringt}{to bring} der Kuh Futter, und der Student schreibt jedes Wort in sein Heft. Die Arbeit scheint leicht, aber der Bauer \voc[erklären]{erklärt}{to explain}: „Die \voc[die Schwierigkeit]{Schwierigkeit}{difficulty} sieht man nicht sofort. Man muß jeden Tag hier bleiben, auch wenn es regnet.“

Hinter dem Stall steht ein \voc[der Arbeiter]{Arbeiter}{worker}. Er \voc[bauen]{baut}{to build} eine neue Tür. Auf einer \voc[die Seite]{Seite}{side} der Tür ist das Holz hell, auf der anderen Seite ist es dunkel. Der Arbeiter zeigt dem Studenten einen \voc[der Fehler]{Fehler}{error}: Ein Nagel sitzt falsch, und deshalb kann die Tür \voc{brechen}{to break}. Der Bauer holt einen Hammer und hilft. Die \voc[die Sonne]{Sonne}{sun} \voc[leuchten]{leuchtet}{to shine} auf das Holz, und der \voc[der Schatten]{Schatten}{shadow} des Arbeiters liegt \voc{breit}{wide} auf der \voc[die Erde]{Erde}{earth}.

Im Haus riecht es nach Kaffee und nach \voc[der Apfelkuchen]{Apfelkuchen}{apple cake}. Die Mutter des Bauern stellt den Kuchen auf den Tisch. Ein \voc[der Knabe]{Knabe}{boy} aus dem Nachbarhaus kommt herein. Er hat einen roten \voc[der Mund]{Mund}{mouth}, kalte Hände und großen Hunger. Die Mutter \voc[schenken]{schenkt}{to give as a present} ihm ein Stück Kuchen. Dem Knaben \voc[schmecken]{schmeckt}{to taste} der Kuchen sehr gut, und sein \voc[der Bauch]{Bauch}{stomach} ist bald voll.

Dann entsteht ein kleines \voc[das Unglück]{Unglück}{misfortune}. Ein zweites Stück Kuchen ist plötzlich weg. Anna sieht auf den Tisch und sagt: „Das Stück ist \voc{gestohlen}{stolen}.“ Niemand will es gewesen sein. Der Knabe sieht in den alten \voc[der Spiegel]{Spiegel}{mirror} neben dem Fenster und sieht Kuchen an seinem Mund. Er wird rot. Die Mutter ist nicht \voc{böse}{bad, guilty}; sie sagt nur: „Ein hungriges Kind macht manchmal ein dummes \voc[das Ding]{Ding}{thing}.“

Der Bauer sagt nichts über \voc[die Strafe]{Strafe}{punishment}. Er gibt dem Knaben Wasser und läßt ihn am Tisch sitzen. Dem Studenten gefällt diese \voc[die Ruhe]{Ruhe}{rest}. Der Knabe sagt leise die \voc[die Wahrheit]{Wahrheit}{truth}: „Ich habe das Stück genommen.“ Danach ist das \voc[das Gewissen]{Gewissen}{conscience} des Knaben wieder ruhig. Die Mutter sagt: „Vielleicht ist das \voc[das Glück]{Glück}{happiness} des Hauses heute größer als das Unglück.“ Der Student schreibt diesen Satz auf.

Am Mittag kommt der \voc[der Arzt]{Arzt}{doctor} aus dem Dorf. Er bringt \voc[die Arznei]{Arznei}{medicine} für die Mutter des Bauern, denn ihr Bein tut weh. Der Arzt sieht auch das Kalb und erklärt dem Studenten: „Bei einem jungen \voc[das Tier]{Tier}{animal} muß man vorsichtig sein. Zu viel kaltes Wasser kann dem Bauch \voc[schaden]{schaden}{to harm}.“ Das Kalb schaut mit einem großen \voc[das Auge]{Auge}{eye} zum Arzt, als ob es jedes Wort versteht.

Nach dem Essen kommt der Lehrer mit vielen Kindern. Die Schule bereitet ein kleines \voc[das Drama]{Drama}{drama} für das Dorf vor. Das Drama heißt „Der \voc[der Löwe]{Löwe}{lion} und der \voc[der Zwerg]{Zwerg}{dwarf}“. Der Hof ist für die Kinder ein guter Ort, denn im Drama gibt es ein Tier, einen Stall, einen Wald und ein kleines Haus. Anna spielt den Zwerg, der Knabe spielt den Löwen. Der Lehrer ist der \voc[der Meister]{Meister}{master} der kleinen Bühne, aber die Kinder lachen so viel, daß die Probe \voc{schwer}{difficult} wird.

Der Lehrer zeigt dem Studenten ein kleines Heft. Darin stehen eine \voc[kurz]{kurze}{short} \voc[die Predigt]{Predigt}{sermon}, ein \voc[der Spruch]{Spruch}{maxim} und ein \voc[das Sprichwort]{Sprichwort}{proverb}; sie gehören zum alten Drama. Der Lehrer sagt: „Das ist kein großes \voc[das Meisterstück]{Meisterstück}{masterpiece}, aber für Kinder ist es gut.“ Im Stück spricht der Löwe von \voc[der Gott]{Gott}{God}, vom \voc[der Teufel]{Teufel}{devil}, von einem \voc[der Fluch]{Fluch}{curse}, von der \voc[die Hölle]{Hölle}{hell} und von einem \voc[der Segen]{Segen}{blessing}. Der Student findet diese alte \voc[die Sprache]{Sprache}{language} schwer, aber interessant.

Die \voc[der Zuschauer]{Zuschauer}{viewer} aus dem Dorf kommen schon vor der Probe. Viele \voc[die Leute]{Leute}{people} stehen am Gartenzaun. Die \voc[die Eltern]{Eltern}{parents} der Kinder bringen Brot, Wasser und eine \voc[die Bratwurst]{Bratwurst}{sausage}. Ein \voc[der Freund]{Freund}{friend} des Bauern bringt eine Lampe. Das Licht der Lampe und das Licht der Sonne fallen auf den Hof; trotzdem bleibt unter dem alten Baum ein dunkler Schatten. Alle warten ruhig, denn die Probe soll nicht in Hast beginnen.

Bei der Probe passiert ein zweites Problem. Der Knabe läuft als Löwe über den Hof, aber er sieht den Kieselstein nicht. Er fällt, und sein Kleid des Löwen wird schmutzig. Anna hilft ihm. Der Arzt sieht sein Bein und sagt: „Es ist nicht schlimm. Ruhe hilft.“ Der Knabe will nicht Ruhe, sondern weiter spielen. Sein \voc[der Wille]{Wille}{will} ist stark, aber sein Bein ist schwach. Der Lehrer kürzt die Szene, und der Löwe darf sitzen.

Dann stellt der Lehrer dem Zwerg ein \voc[das Rätsel]{Rätsel}{riddle}: „Was hat ein \voc[das Angesicht]{Angesicht}{face} ohne Mund und ein Auge ohne Seele?“ Anna denkt lange nach. Der Lehrer zeigt auf den Spiegel. Alle Kinder lachen, denn im Spiegel sieht Anna ihr eigenes Angesicht. Der Arzt sagt: „Ein Spiegel hat keine \voc[die Seele]{Seele}{soul}; ein Kind aber hat eine Seele, und man muß sie schützen wie mit einem \voc[der Schild]{Schild}{shield}.“

Der Student schreibt nun nicht nur über die Arbeit des Bauern, sondern auch über das \voc[das Leben]{Leben}{life} auf dem Hof. Er sieht \voc[die Tugend]{Tugend}{virtue} in kleinen Dingen: im Fleiß des Bauern, in der \voc[die Geduld]{Geduld}{patience} der Mutter, in der Hilfe des Arbeiters und in der Ehrlichkeit des Knaben. Er schreibt auch über \voc[die Ehre]{Ehre}{honor}, denn der Knabe hat die Wahrheit gesagt, obwohl es ihm \voc{unangenehm}{unpleasant} war.

Am Nachmittag kommt ein Brief aus der Schule. Der Lehrer \voc[senden]{sendet}{to send} dem Bauern eine Nachricht: Die Kinder sollen nächste Woche wieder auf dem Hof proben. Anna liest den Brief und zeigt ihn der Mutter. Der Bauer erklärt dem Studenten: „Für die Schule ist der Hof ein gutes Buch. Hier sieht man Arbeit, Tier, Erde, Sonne, Fehler und \voc[die Besserung]{Besserung}{improvement}.“ Der Student schreibt jedes Wort sorgfältig auf.

Dann fragt Anna den Studenten nach dem Sinn seines Berichts. Er sagt: „Ich schreibe über \voc[die Selbsterkenntnis]{Selbsterkenntnis}{self-knowledge}.“ Anna versteht das Wort nicht sofort. Der Student erklärt: „Der Bauer lernt durch seine Arbeit etwas über sich selbst. Der Knabe lernt durch den Kuchen etwas über sein Gewissen. Und ich lerne durch meinen Bericht, daß das \voc[das Lernen]{Lernen}{learning} nicht nur in Büchern geschieht.“ Anna nickt. Das ist keine Predigt, sondern ein Satz aus dem Tag.

Gegen Abend wird der Himmel dunkel. Eine \voc[die Fliege]{Fliege}{fly} auf dem Fenster fliegt noch einmal über den Tisch. Der Arzt sagt: „Diese Fliege bleibt immer bei dem Kuchen.“ Die Mutter \voc[bedecken]{bedeckt}{to cover} den Kuchen mit einem Teller, aber sie bedeckt nicht das Lachen der Kinder. Der Bauer holt die Kuh in den Stall. Das Kalb läuft fast durch die Tür, aber Anna hält es zurück. So kann es Anna nicht \voc[entgehen]{entgehen}{to escape}.

Der Student sitzt am Tisch und schreibt seinen Bericht zu Ende. Er schreibt nicht von \voc[das Wohlleben]{Wohlleben}{well-being}, denn auf dem Hof ist das Leben nicht bequem. Er schreibt von Arbeit, \voc[die Zeit]{Zeit}{time} und Geduld. Er schreibt: „Die Zeit des Bauern ist anders als die Zeit der Stadt. Auf dem Hof ist jeder Morgen eine Arbeit, jede Arbeit eine Übung, und jede Übung eine \voc[die Wiederholung]{Wiederholung}{repetition}.“ Dieses Wort gefällt ihm, weil es zum Tag paßt.

Später zeigt die Mutter dem Studenten ein altes \voc[das Werk]{Werk}{work} ihres Vaters. In dem Buch steht ein Abschnitt über \voc[die Weisheit]{Weisheit}{wisdom}, über den \voc[der Zorn]{Zorn}{wrath} und über die Ruhe. Der Student liest langsam. Der Text sagt nicht, daß der Zorn gut ist; er sagt nur, daß man Zeit braucht, wenn der Zorn im Herzen ist. Anna sagt: „Das klingt wie ein Sprichwort.“ Der Student sagt: „Ja, aber heute verstehe ich es besser.“

Als der Student geht, schenkt ihm die Mutter ein kleines Stück Apfelkuchen. Der Knabe sieht es und lacht. „Dieses Stück ist nicht gestohlen“, sagt er. Der Bauer bringt den Studenten bis zum Weg. \voc{Hoffentlich}{hopefully} wird der Bericht gut, denkt der Student. Die Sonne ist weg, aber ein letzter heller Streifen liegt noch auf der Erde. Der Student denkt: „Das war ein schwerer, aber guter Tag.“ Dann geht er zurück in die Stadt und sendet am Abend seinen Bericht an die Schule.

Am nächsten Morgen liest der Lehrer den Bericht. Er findet alles klar und sagt: „Das ist wohl kein Meisterstück, aber es ist ganz gut.“ Die Kinder lachen. Niemand will die Fliege \voc{töten}{to kill}, denn sie gehört nun fast zum Drama. Der Knabe \voc[lieben]{liebt}{to love} den Apfelkuchen, Anna liebt die Bühne, und der Bauer liebt seinen Hof. So bleibt die Arbeit schwer, aber das Leben auf dem Hof wird durch Wiederholung, Geduld und Freundschaft jeden Tag ein wenig besser.

\end{readersection}


\begin{readersection}{Der Bericht aus dem Stadtmuseum}

Am Montag schreibt die Zeitung über eine neue Ausstellung im Stadtmuseum. Die Ausstellung heißt „Arbeit, \voc[die Armut]{Armut}{poverty} und \voc[der Reichtum]{Reichtum}{wealth} in der alten Stadt“. Sie zeigt nicht nur schöne Dinge aus alten Häusern, sondern auch das \voc[das Herzeleid]{Herzeleid}{suffering} vieler Menschen. Ein \voc[arm]{armer}{poor} Arbeiter, ein reicher \voc[der Herr]{Herr}{Lord, Master, Mr.}, eine \voc[ander]{andere}{other} Familie und ein \voc[still]{stilles}{quiet} Kind kommen in den Bildern vor.

Der Leiter der Ausstellung ist ein \voc[der Historiker]{Historiker}{historian}. Er ist kein berühmter Mann, aber er arbeitet \voc{täglich}{daily} im Museum und liest \voc{gern}{gladly} alte Briefe. In seinem Zimmer gibt es zwei hohe \voc[das Fenster]{Fenster}{window}, ein \voc[dunkel]{dunkles}{dark} Regal und viele Bücher. Auf dem Tisch liegt ein Blatt mit einem \voc[finster]{finsteren}{gloomy} \voc[der Gedanke]{Gedanken}{thought}: „Auch kleine Dinge können eine große Geschichte haben.“

Im ersten Saal steht ein altes Bild von einer \voc[die Kirche]{Kirche}{church}. Über dem \voc[das Dach]{Dach}{roof} der Kirche liegen graue \voc[die Wolke]{Wolken}{cloud}, und unter dem Dach sitzen viele Leute. Der \voc[der Regen]{Regen}{rain} fällt auf die Straße. Eine Frau trägt \voc[sauer]{saure}{sour} \voc[die Frucht]{Frucht}{fruit} in einem Korb, ein Kind hält \voc[süß]{süßes}{sweet} \voc[das Obst]{Obst}{fruit}, und ein alter Mann trinkt \voc[bitter]{bitteren}{bitter} Kaffee. Der Historiker erklärt den Besuchern: „Das Bild ist nicht nur schön. Es zeigt, wie schwer das Leben damals war.“

Neben dem Bild steht eine Vitrine. Darin liegt ein kleiner Apparat von einem \voc[der Erfinder]{Erfinder}{inventor}. Der Erfinder wollte mit dem Apparat Wasser aus dem Keller holen, aber das Ding war zu \voc[leicht]{leicht}{easy} gebaut und ging bald kaputt. Ein Besucher lacht. Der Historiker sagt ruhig: „Bitte kein \voc[das Geschwätz]{Geschwätz}{gossip}. Ein schlechter Apparat kann trotzdem ein wichtiger Anfang sein. Manchmal \voc[hindern]{hindert}{to hinder} ein Fehler die Arbeit nicht, sondern hilft dem Lernen.“

Im zweiten Saal hängt ein Bild von einem \voc[mager]{mageren}{thin} \voc[der Hund]{Hund}{dog}. Der Hund sitzt bei einem \voc[der Bär]{Bären}{bear} aus Holz. Ein Mädchen fragt: „Warum sieht der Hund so traurig aus?“ Der Historiker antwortet: „Er gehört zu einer Geschichte über den \voc[der Tod]{Tod}{death} seines Besitzers.“ Auf dem Bild hat der Hund eine kleine \voc[die Wunde]{Wunde}{wound} am \voc[der Kopf]{Kopf}{head}; die Wunde \voc[bluten]{blutet}{to bleed} nicht mehr, aber sie sieht noch rot aus. Ein Kind wird \voc{stumm}{mute} und sagt nichts.

Im dritten Saal steht eine alte Waschschüssel. Daneben steht: „Mit diesem \voc[das Wasser]{Wasser}{water} \voc[waschen]{wäscht}{to wash} die Mutter jeden Morgen die Hände der Kinder.“ Das Wasser war kalt, aber es war eine \voc[die Gabe]{Gabe}{gift}, denn nicht jedes Haus hatte sauberes Wasser. Auf einem anderen Schild steht: „Es war \voc{verboten}{forbidden}, das Wasser des Herrn zu nehmen.“ Der Historiker sagt: „Solche Regeln zeigen, wer in der Stadt Macht hatte.“

An der Wand daneben liest man einen Text über eine alte Familie. Der Vater wollte das Haus \voc[regieren]{regieren}{to rule} wie ein kleines Reich. Er war \voc{ständig}{constant} im Hof und sah alles. Sein Sohn wollte eigene Bilder \voc[schaffen]{schaffen}{to create}, aber der Vater sagte, Kunst \voc[verderben]{verdirbt}{to spoil} die Arbeit. Die Mutter schrieb später: „Ein \voc[schlecht]{schlechtes}{bad} Wort kann ein ganzes Jahr dunkel machen.“ Der Historiker findet diesen Satz traurig, aber wahr.

Ein Junge läuft zu schnell durch den Saal. Die Aufsicht sagt: „Langsam bitte!“ Der Junge geht \voc{rückwärts}{backwards}, stößt fast gegen eine Lampe und lacht. Seine Schwester sagt: „Das ist \voc{gefährlich}{dangerous}.“ Sie zeigt auf ein kleines Schild: „Nicht berühren.“ Auf dem Schild sitzt eine \voc[die Mücke]{Mücke}{gnat}. Der Junge will sie \voc[töten]{töten}{to kill}, aber die Schwester sagt: „Laß sie. Sie gehört jetzt auch zur Ausstellung.“

In einer Ecke liegt noch ein altes \voc[das Werk]{Werk}{work} über Pflanzen. Auf der ersten Seite sieht man die \voc[die Wurzel]{Wurzel}{root} eines Baumes. Daneben steht ein \voc[übel]{übles}{evil} Wort über den \voc[der Teufel]{Teufel}{devil}; der Historiker sagt aber, dass solche Wörter nur zur alten Sprache der Ausstellung gehören. \voc{Vielleicht}{perhaps} bleiben einige Besucher \voc{lange}{for a long time} vor diesem Buch stehen, aber \voc[all]{alle}{all, every} verstehen es nicht sofort.

Am Ende der Ausstellung steht ein kleines Theater. Dort \voc[tanzen]{tanzen}{to dance} Kinder in alten Kleidern. Ein Kind spielt einen Herrn, ein anderes spielt ein altes \voc[das Weib]{Weib}{wife}. Der Historiker sagt den Besuchern, dass dieses Wort heute alt und hart klingt. Im Theatertext aber steht es so, und deshalb bleibt es im Heft. Ein \voc[melancholisch]{melancholischer}{melancholy} Musiker sitzt daneben und spielt leise. Seine Musik macht die Szene nicht fröhlich, aber \voc[sanft]{sanft}{soft}.

Nach der Vorstellung fragt eine Besucherin: „\voc[glauben]{Glauben}{to believe} Sie, dass solche alten Dinge heute noch wichtig sind?“ Der Historiker sieht auf die Vitrine, auf das Bild vom Hund, auf die Kirche und auf die Kinder. Dann sagt er: „Ja. Ein Museum zeigt nicht nur alte Dinge. Es zeigt das \voc[das Gewissen]{Gewissen}{conscience} einer Stadt.“ \voc{Endlich}{finally} ist der Saal still. Draußen regnet es noch immer, aber im Museum bleibt das Licht warm.

\end{readersection}

\begin{readersection}{Der Bericht vom Hafen}

Lina \voc[wohnen]{wohnt}{to live, reside} \voc{seit}{since} \voc[das Jahr]{drei Jahren}{year} in einer kleinen \voc[die Stadt]{Stadt}{city} am Fluß. Drei Jahre \voc{lang}{for} geht sie fast jeden \voc[der Morgen]{Morgen}{morning} \voc{durch}{through} die leeren Straßen zum Hafen. Heute steht sie \voc{um}{at} \voc{fünf}{five} \voc[die Uhr]{Uhr}{clock} auf und ist um \voc{sechs}{six} Uhr am Kai. Sie will einen Bericht über den alten Hafen schreiben, denn die Stadt will ihn verkaufen. Linas \voc[der Beruf]{Beruf}{occupation, job} ist der Journalismus; ihre \voc[die Absicht]{Absicht}{intention} ist einfach: Sie will zeigen, was der Hafen für die Menschen bedeutet.

Am Kai \voc[treffen]{trifft}{to hit} Lina Herrn Brandt, einen alten \voc[der Erzähler]{Erzähler}{narrator}. Er hat früher auf einem \voc[das Schiff]{Schiff}{ship} gearbeitet. Neben dem Schiff steht noch der hohe \voc[der Mast]{Mast}{mast}, und an dem Mast hängt ein altes Seil. \voc{Unter}{under} dem Schiff liegt nasses \voc[das Stroh]{Stroh}{straw}; im Wasser \voc[treiben]{treiben}{to drive} ein Brett und ein Stück Papier. Der \voc[der Strom]{Strom}{stream} ist heute schnell und fast \voc{reißend}{rushing}. Herr Brandt sagt: „\voc{Vor}{ago} vielen Jahren war hier alles \voc{lebendig}{lively}. Jetzt ist der Hafen \voc{traurig}{sad}.“

Lina stellt die erste \voc[die Frage]{Frage}{question}: „\voc[was für]{Was für}{what kind of} ein Ort war der Hafen früher?“ Herr Brandt antwortet: „Ein Ort der Arbeit, der \voc[die Hoffnung]{Hoffnung}{hope}, der \voc[der Glaube]{Glaube}{faith} und der \voc[die Armut]{Armut}{poverty}. Die Arbeiter kamen \voc{aus}{from} vielen Häusern, oft \voc{ohne}{without} Geld, aber nicht ohne \voc[das Herz]{Herz}{heart}. Sie arbeiteten \voc{für}{for} ihre Kinder und \voc{wider}{against} das \voc[das Unglück]{Unglück}{misfortune}. \voc[Jeder]{Jeder}{every, each} Mann kannte hier jeden \voc[der Nagel]{Nagel}{nail}, jede \voc[die Tür]{Tür}{door} und jeden \voc[der Schlüssel]{Schlüssel}{key}.“

\voc{Gegen}{towards} \voc[der Mittag]{Mittag}{noon} geht Lina \voc{mit}{with} Herrn Brandt in ein kleines Büro. \voc{An}{at} der Wand hängt ein Bild vom Hafen, \voc{auf}{on} dem \voc[der Tisch]{Tisch}{table} liegen alte Briefe, \voc{neben}{next to} dem Tisch steht ein Stuhl, \voc{hinter}{behind} dem Stuhl steht ein Korb, und \voc{zwischen}{between} dem Korb und der Wand schläft ein alter Hund. Herr Brandt setzt sich Lina \voc{gegenüber}{opposite}. \voc{Außer}{except for} dem Hund ist niemand im Zimmer. Lina lacht \voc{gern}{gladly}, aber heute bleibt sie ernst.

Sie liest einen Brief, der \voc{von}{of, from} einem \voc[der Christ]{Christen}{Christian} aus \voc[das Frankreich]{Frankreich}{France} kommt. Er schreibt: „In Frankreich \voc[leben]{lebt}{to live} man vielleicht gut, aber am alten Hafen lebt man wahrer.“ Lina findet den Satz seltsam. Herr Brandt sagt: „Das ist eine \voc[die Weise]{Weise}{way}, über Heimat zu \voc[reden]{reden}{to speak}. Der Mann hatte eine eigene \voc[die Weltanschauung]{Weltanschauung}{worldview}. Er hat den Hafen nicht nur gesehen; er hat die \voc[die Wirklichkeit]{Wirklichkeit}{reality} des Hafens verstanden.“

Vor dem Fenster \voc[stehen]{steht}{to stand} ein kleines \voc[das Tier]{Tier}{animal}: eine \voc[die Ente]{Ente}{duck}. Weiter draußen läuft eine \voc[die Gans]{Gans}{goose} über den Kai. \voc{Hinter}{behind} einem \voc[der Busch]{Busch}{bush} \voc[schleichen]{schleicht}{to creep} ein \voc[der Fuchs]{Fuchs}{fox}; er sieht die Ente, aber er \voc[bekommen]{bekommt}{to receive} sie nicht. Die Ente läuft \voc{über}{over} eine kleine \voc[die Brücke]{Brücke}{bridge} und \voc{weg}{away}. Herr Brandt sagt: „Ein Fuchs geht \voc{nie}{never} leicht in eine \voc[die Falle]{Falle}{trap}.“ Lina schreibt es auf, obwohl es wie ein \voc[das Sprichwort]{Sprichwort}{proverb} klingt.

Dann kommen zwei Kinder \voc{bei}{at} Herrn Brandt vorbei. Sie tragen ein kleines \voc[das Geschenk]{Geschenk}{present}; es ist \voc[die Extrapost]{per Extrapost}{special delivery} gekommen. Das Geschenk ist ein Buch über \voc[die Ebbe]{Ebbe}{ebb}, \voc[die Flut]{Flut}{flow}, \voc[der Himmel]{Himmel}{sky} und \voc[das Wetter]{Wetter}{weather}. Die Kinder öffnen das Buch und lesen: „Ebbe und Flut \voc[warten auf]{warten auf}{to wait for} niemanden.“ Herr Brandt nickt: „Auch der Hafen wartet nicht. Wenn die Stadt ihn verkauft, ist ein Jahr später vieles anders.“

Im Hof vor dem Büro gibt es ein kleines \voc[das Feld]{Feld}{field}; dort liegt im Frühling \voc[die Saat]{Saat}{seed}. Neben dem Feld steht eine alte \voc[die Windmühle]{Windmühle}{windmill}. Lina geht \voc{um}{around} die Windmühle und sieht, daß ihr Holz alt ist. Herr Brandt sagt: „Das \voc[das Haar]{Haar}{hair} des Menschen wird weiß, und das Holz wird grau. Alles \voc[wachsen]{wächst}{to grow}, wird alt und bleibt \voc[ein wenig]{ein wenig}{somewhat} fremd.“ Lina fragt: „Ist das nicht \voc[das Elend]{Elend}{misery}?“ Herr Brandt sagt: „Nein. Das ist Leben.“

\voc{Nach}{after} dem Spaziergang \voc[frühstücken]{frühstückt}{to breakfast} Lina spät im Hafenlokal. Sie \voc[essen]{ißt}{to eat} Brot und Käse, und Herr Brandt trinkt \voc[der Wein]{Wein}{wine}. An einem Tisch sitzt ein \voc[der Richter]{Richter}{judge} mit einem jungen Paar. Der Richter spricht über Geld, \voc[die Liebe]{Liebe}{love} und \voc[die Gegenliebe]{Gegenliebe}{reciprocal love}. Das Paar will bald \voc[heiraten]{heiraten}{to marry}, aber die Familien haben \voc[die Furcht]{Furcht}{fear}. Lina hört nicht alles, doch sie hört den Satz: „\voc[die Leidenschaft]{Leidenschaft}{passion} ist groß, aber Geduld ist größer.“

Herr Brandt zeigt ihr ein altes Heft. Darin steht: „Der \voc[der Stolz]{Stolz}{pride} frühstückt mit dem \voc[der Überfluß]{Überfluß}{abundance}, ißt zu Mittag mit der Armut und ißt am \voc[der Abend]{Abend}{evening} mit der \voc[die Schande]{Schande}{shame}.“ Lina findet den Satz \voc{scharf}{sharp} wie einen \voc[der Zahn]{Zahn}{tooth}. Herr Brandt sagt: „Alte Sätze sind oft hart. Man muß sie nicht für die ganze Wahrheit \voc[halten für]{halten}{to consider}, aber man kann sie als \voc[die Beobachtung]{Beobachtung}{observation} lesen.“

Im Lokal hängt ein Schild: „\voc{Während}{during} des Sturms bleibt der Kai geschlossen.“ \voc{Trotz}{in spite of} des Regens sitzen einige Leute draußen. \voc{Wegen}{because of} einer \voc[die Krankheit]{Krankheit}{sickness} bleibt der Wirt heute zu Hause; \voc{statt}{instead of} des Wirts arbeitet seine Tochter. \voc{Anstatt}{instead of} eines langen Menüs gibt es nur Suppe. \voc{Außerhalb}{outside of} des Lokals ist der Wind laut, aber im Zimmer brennt \voc[das Feuer]{Feuer}{fire}, und über dem Herd hängt \voc[der Rauch]{Rauch}{smoke}.

Lina fragt die Tochter: „Warum arbeiten Sie heute hier?“ Die Tochter antwortet: „\voc[um . . . willen]{Um meines Vaters willen}{for the sake of} \voc[tun]{tue}{to do} ich es. Er hat Schmerzen und \voc[leiden]{leidet}{to suffer} sehr.“ Dann lächelt sie: „Ich bin nicht \voc{faul}{lazy}; \voc[meistens]{meistens}{mostly} bin ich sogar \voc{fleißig}{hard-working}.“ Herr Brandt \voc[lachen]{lacht}{to laugh}. „Das stimmt“, sagt er. „Sie \voc[machen]{macht}{to do, make} die beste Suppe der Stadt.“ Die Tochter wird rot, und Lina bemerkt ihre gute \voc[die Laune]{Laune}{mood}.

Später \voc[besuchen]{besucht}{to visit} Lina den Hafenmeister. Sein Büro liegt \voc{in}{in} einem niedrigen Haus vor der Brücke. Er spricht meistens ruhig, aber heute ist er ärgerlich. „Die Stadt redet über den Hafen, aber sie kennt ihn kaum. \voc[folgend]{Folgendes}{following} ist wichtig: Der Hafen ist kein leerer Ort. Hier arbeiten Menschen, hier gibt es \voc[die Tat]{Taten}{deed}, hier gibt es Fehler, hier gibt es \voc[die Tugend]{Tugend}{virtue}.“ Lina fragt nach dem \voc[der Begriff]{Begriff}{concept} „Wert“. Der Hafenmeister sagt: „Der \voc[der Wert]{Wert}{value} eines Ortes ist nicht nur \voc[das Geld]{Geld}{money}.“

Im selben Büro liegt eine \voc[die Rede]{Rede}{speech} des Bürgermeisters. Die Rede \voc[schildern]{schildert}{to portray} den Hafen als Problem. Sie spricht von Kosten, von Gefahr und von wenig Nutzen. Herr Brandt \voc[mißtrauen]{mißtraut}{to mistrust} dieser Rede. „Der Bürgermeister kennt den Hafen nicht“, sagt er. „Er sieht nur Zahlen.“ Lina schreibt: „Der \voc[der Mensch]{Mensch}{person} lebt nicht von Zahlen allein.“ Dann denkt sie: Das klingt wieder wie ein Sprichwort, \voc{doch}{nevertheless} vielleicht braucht ein Bericht manchmal solche Sätze.

Am Nachmittag kommt ein alter Mann herein. Er ist nicht reich, aber er hat Würde. Er sagt: „Ich war hier Arbeiter. \voc{Vor allem}{above all} habe ich gelernt, daß Arbeit ohne Freundschaft arm ist.“ Dann erzählt er von einem Freund, der \voc{irgendwo}{somewhere} im Süden lebt. „Wir hatten \voc[verschieden]{verschiedene}{different} Berufe, aber denselben Hafen. \voc{Deshalb}{for that reason} schreiben wir einander noch.“ Lina fragt: „\voc[fürchten]{Fürchten}{to fear} Sie die Veränderung?“ Der Mann antwortet: „Ich fürchte nicht Veränderung, sondern Vergessen.“

Dann beginnt ein kurzes \voc[das Gespräch]{Gespräch}{conversation} über das Theater am Hafen. Früher spielte dort eine kleine Gruppe jedes Jahr ein \voc[das Drama]{Drama}{drama} über Schiffe, Sturm und Arbeit. Herr Brandt sagt: „Ich spielte einmal einen Matrosen. Ich hatte einen Hut auf dem \voc[der Kopf]{Kopf}{head} und stand neben dem Mast. Das Publikum war klein, aber freundlich.“ Lina fragt: „Treffen sich die Schauspieler noch?“ Herr Brandt antwortet: „Ja, manchmal treffen wir uns hier.“

Gegen Abend geht Lina noch einmal den Kai \voc{entlang}{along}. Sie sieht die leeren Schiffe, die nassen Seile und einen Jungen, der \voc{fast}{almost} ins Wasser fällt. Seine Mutter \voc[schlagen]{schlägt}{to hit} ihn nicht; sie hält ihn nur fest und sagt: „Ich rette dich.“ Der Junge lacht. Dann will er den Hund \voc[küssen]{küssen}{to kiss}, aber der Hund läuft weg. Herr Brandt sagt: „Kinder tun seltsame Dinge.“

Als Lina zurückgeht, spricht sie mit Herrn Brandt über \voc[das Schlafen]{Schlafen}{sleeping}. „Nach so einem Tag schläft man gut“, sagt er. „Aber wenn ein Ort verschwindet, kommt nachts der \voc[der Jammer]{Jammer}{sorrow}.“ Lina fragt: „Ist das Ihre Wahrheit?“ Er antwortet: „Nur meine Weise zu reden.“ Sie gehen bis zur Tür des Büros. Herr Brandt bleibt stehen und sagt: „Ich habe \voc[wenig]{wenig}{little} Hoffnung, aber ich habe noch Hoffnung.“

Am nächsten Morgen erscheint Linas Bericht in der Zeitung. Der Titel lautet: „Der Hafen und seine Menschen“. Der Bericht ist nicht \voc{ganz}{whole, entire} lang, aber er ist klar. Er erzählt von Armut und Überfluß, von Liebe und Gegenliebe, von Arbeit und \voc[die Faulheit]{Faulheit}{laziness}, von Furcht und Hoffnung. Er zeigt, daß der Hafen mehr ist als ein Grundstück. Er ist ein Ort, an dem Menschen wohnen, arbeiten, reden, warten, leiden und lachen. Und er endet mit Herrn Brandts Satz: „Wenn ein Ort lebendig war, bleibt er nicht einfach weg.“

\end{readersection}


\begin{readersection}{Fräulein Meier und das Schäflein}

Am Morgen steht Fräulein Meier in der \voc[die Küche]{Küche}{kitchen}. Auf dem Tisch liegt ein Brief, und neben dem Brief steht ein alter \voc[der Kalender]{Kalender}{calendar}. Fräulein Meier sieht auf den Kalender und fragt: „\voc{Wann}{when} kommt der Briefträger? \voc{Warum}{why} ist er heute so spät? \voc{Was}{what} bringt er mir? \voc{Wo}{where} ist meine Katze?“ \voc{Dort}{there}, unter dem Tisch, liegt die Katze und schläft.

Da hört Fräulein Meier eine Stimme vor der Tür. „\voc[wer]{Wer}{who} ist da?“ ruft sie. Der Briefträger antwortet: „Ich bin es. Öffnen Sie die Tür!“ Fräulein Meier öffnet die Tür. Vor der Tür steht der Briefträger mit einem großen \voc[der Sack]{Sack}{sack}. Hinter ihm steht ein \voc[der Bauer]{Bauer}{farmer}, und neben dem Bauern steht ein kleines \voc[das Schäflein]{Schäflein}{little sheep}. Das Schäflein ist nass, denn draußen fällt \voc[der Regen]{Regen}{rain}.

„Wem \voc[gehören]{gehört}{to belong to} dieses Schäflein?“ fragt Fräulein Meier. Der Bauer antwortet: „Es gehört mir. Aber heute habe ich keine Zeit. Ich bin ein \voc[der Arbeiter]{Arbeiter}{worker}, und ich muss in die Stadt gehen.“ Der Briefträger sagt: „Der Bauer \voc[kennen]{kennt}{to know} mich, und ich kenne Sie. Darum bringe ich das Schäflein zu Ihnen.“ Fräulein Meier sieht den Bauern, den Briefträger und das Schäflein an. „Ich kenne Sie beide“, sagt sie, „aber dieses Schäflein kenne ich nicht.“

Das Schäflein hat viel \voc[die Wolle]{Wolle}{wool} und einen stillen \voc[der Blick]{Blick}{glance}. „\voc{Wie}{how, like} \voc[heißen]{heißt}{to be called} es?“ fragt Fräulein Meier. Der Bauer sagt: „Es heißt noch nichts.“ Der Briefträger fragt: „Wie soll es heißen?“ Fräulein Meier denkt einen Augenblick und sagt: „Es heißt vielleicht Hoffnung.“ Der Briefträger lacht. „Das ist ein guter Name. Ich \voc[glauben]{glaube}{to believe}, dass Hoffnung auch bei Regen gut ist.“

Das Schäflein läuft in das Haus und dann direkt in die Küche. Die Katze springt auf den Stuhl. „\voc[meiden]{Meiden}{to avoid} Sie die Katze!“ ruft Fräulein Meier. „Meide die Katze, kleines Schäflein!“ Aber das Schäflein hört nicht. Es läuft unter den Tisch, stößt an ein \voc[das Bein]{Bein}{leg} des Tisches und fällt fast in ein kleines \voc[das Bad]{Bad}{bath}, das Fräulein Meier für die Katze vorbereitet hat. „\voc[nehmen]{Nehmen}{to take} Sie das Bad weg!“ ruft sie dem Briefträger zu. Der Briefträger nimmt das Bad und stellt es neben die Tür.

Der Bauer sieht das Schäflein und sagt: „Es ist \voc{arm}{poor} und klein, aber es ist \voc{fröhlich}{joyous}. Es hat \voc{nichts}{nothing} Böses getan.“ Fräulein Meier sagt: „Ein armes Tier braucht Brot, Wasser und \voc[die Liebe]{Liebe}{love}.“ Der Briefträger antwortet: „Und \voc{genug}{enough} Ruhe.“ Der Bauer nickt. „Ja, Ruhe ist gut. Aber ein \voc[faul]{fauler}{lazy} Bauer hat nie genug Ruhe.“ Fräulein Meier lacht: „Sind Sie faul, lieber Bauer?“ Der Bauer wird rot und sagt: „Nein, ich arbeite \voc{immer}{always}.“

In diesem Augenblick kommt der alte Historiker. Er trägt ein Buch und liest laut einen \voc[der Satz]{Satz}{sentence}: „Der \voc[die Tugend]{Tugend}{virtue} ist kein \voc[das Ziel]{Ziel}{goal} \voc{zu}{too} \voc{hoch}{high}.“ Fräulein Meier sagt: „Heute ist mein Ziel nicht hoch. Ich will nur dieses Schäflein trocken machen.“ Der Historiker sieht das Tier und fragt: „Warum steht ein Schäflein in der Küche?“ Der Briefträger antwortet: „Weil der Bauer keine Zeit hat und weil Fräulein Meier ein gutes \voc[das Herz]{Herz}{heart} hat.“ Der Historiker sagt: „Das ist ein klarer Satz.“

Dann sieht der Historiker den Sack. „Was ist in dem Sack?“ fragt er. Der Briefträger öffnet den Sack. Darin liegen Wolle, Brot und ein kleines Buch. Auf dem Buch steht: \voc[das Gebet]{Gebete}{prayer} für Kinder. Der Historiker sagt: „Ein Gebet gehört nicht in einen Sack.“ Fräulein Meier antwortet: „Ein gutes Gebet kann überall liegen.“ Der Bauer sagt: „Meine Mutter gab mir das Buch. Sie sagte: \voc{Ohne}{without} Gebet wird das \voc[das Leben]{Leben}{life} dunkel.“ Fräulein Meier findet diesen Satz schön.

Der Briefträger stellt eine Frage. „\voc[stellen]{Stelle}{to ask, pose} ich jetzt eine dumme \voc[die Frage]{Frage}{question}?“ fragt er. „Warum fürchtet die Katze das Schäflein?“ Fräulein Meier sagt: „Die Katze \voc[fürchten]{fürchtet}{to fear} es nicht. Aber sie kennt es nicht.“ Der Historiker sagt: „\voc[die Furcht]{Furcht}{fear} kommt oft aus Unwissenheit.“ Der Bauer nickt. „Das stimmt. Auch ein \voc[der Schuldige]{Schuldiger}{guilty one} fürchtet manchmal ein kleines Geräusch.“ Der Historiker wird ernst und sagt: „Und ein \voc[der Angeklagte]{Angeklagter}{accused} fürchtet einen Richter.“ Fräulein Meier antwortet: „Heute haben wir keinen Richter, nur ein Schäflein.“

Plötzlich läuft die Katze aus der Küche. Das Schäflein läuft hinter ihr \voc{her}{hither}. Der Hund sieht das Schäflein und bellt. Das Schäflein \voc[erschrecken]{erschrickt}{to frighten}. Fräulein Meier ruft: „\voc[lassen]{Lassen}{to let, leave} Sie den Hund nicht los, aber \voc[verjagen]{verjagen}{to chase} Sie das Schäflein nicht!“ Der Bauer will das Schäflein verjagen, aber Fräulein Meier sagt: „Nein, verjagen Sie es nicht. Es ist noch klein.“ Der Briefträger sagt: „\voc{Niemand}{no one} verjagt heute ein Schäflein aus diesem Haus.“

Jetzt steht ein \voc[der Bettler]{Bettler}{beggar} vor der Tür. Er ist nass und sieht müde aus. Fräulein Meier fragt: „Was wollen Sie?“ Der Bettler sagt: „Ich habe Hunger.“ Der Briefträger gibt ihm Brot. Der Historiker sagt: „Ein \voc[der Geber]{Geber}{giver} hat ein gutes Herz.“ Der Bettler nimmt das Brot und sagt: „Danke.“ Fräulein Meier sagt: „Ohne Brot ist der Mensch traurig, aber ohne Freundschaft ist er noch ärmer.“

Ein \voc[der Lehrer]{Lehrer}{teacher} aus der Stadt kommt auch ins Haus. Er hat eine kleine Tasche und viele Fragen. „Die \voc[die Leute]{Leute}{people} sagen, dass Fräulein Meier heute ein Schäflein hat“, sagt er. „Ist das wahr?“ Fräulein Meier antwortet: „Ja, die Leute reden viel, und manchmal reden sie richtig.“ Der Lehrer sieht das Schäflein. „Seine Beine sind klein, aber gut.“ Dann \voc[schauen auf]{schaut er auf}{to look at} die Wolle und sagt: „Auch die Wolle ist gut.“ Der Bauer sagt: „Der Wolle \voc{wegen}{because of} ist das Schäflein wichtig, aber wegen des Regens ist es heute hier.“ Fräulein Meier nickt. „Wegen des Regens bleibt heute jeder im Haus.“

Der Lehrer setzt sich an den Tisch. „Heute lernen wir deutsche Fragen“, sagt er. „Wie heißt das Tier? Wo steht es? Wann kommt der Bauer wieder? Warum bleibt der Briefträger hier? Wen liebt das Schäflein am meisten? Wem gehört der Sack?“ Fräulein Meier sagt: „Bleiben Sie beim Schäflein, Herr Briefträger!“ Der Briefträger antwortet: „Gern.“ Der Lehrer sagt: „Sehr gut. Fragen helfen beim Lernen der \voc[die Sprache]{Sprache}{language}.“ Der Historiker sagt: „Aber \voc[das Raten]{Raten}{advice} hilft manchmal mehr als Fragen.“ Fräulein Meier antwortet: „Gutes Raten ist kurz, klare Fragen sind besser.“

Der Briefträger sagt: „Ich \voc[gern haben]{habe dieses Schäflein gern}{to like}, weil es so still ist.“ Aber in diesem Augenblick springt das Schäflein auf den Stuhl und stößt den Kalender auf die \voc[die Erde]{Erde}{earth}. Der Kalender fällt in das Bad. „Ach!“ ruft Fräulein Meier. „Jetzt ist mein Kalender nass.“ Der Briefträger sagt: „Das Schäflein hat vielleicht auch den Kalender gern.“ Fräulein Meier lacht \voc{doch}{just, nevertheless}. „\voc[scherzen]{Scherzen}{to joke} Sie nicht zu viel!“ sagt sie. „Scherzen Sie lieber mit dem Historiker.“

Der Historiker hat \voc[der Kummer]{Kummer}{grief, sorrow}, denn sein Buch liegt auch auf dem nassen Tisch. „Das ist ein großer \voc[der Schmerz]{Schmerz}{pain} für einen alten Historiker“, sagt er. Fräulein Meier antwortet: „Lieber Historiker, nehmen Sie Ihren Kummer nicht so ernst. Das Buch ist nur ein wenig feucht.“ Der Lehrer sagt: „\voc{Nur}{only, just} ein wenig? Das Wasser ist überall.“ Fräulein Meier sagt: „\voc[erst]{Erst}{first} trocknen wir den Tisch, \voc{dann}{then} trinken wir Wein.“

Draußen hängt eine dunkle \voc[die Wolke]{Wolke}{cloud} über dem Garten. Der Historiker sagt: „Vielleicht sollte man jetzt ein zweites Gebet lesen.“ Fräulein Meier ruft: „Lesen Sie keinen schweren Satz! Lesen Sie etwas Fröhliches.“ Der Historiker liest: „Die \voc[die Tat]{Tat}{deed} ist besser als viele Worte.“ Der Bauer sagt: „Das glaube ich.“ Der Briefträger sagt: „Dann tun wir etwas. Wir trocknen den Tisch, und danach essen wir.“

Fräulein Meier bringt Brot, Obst, Wein und Suppe. Der Historiker sieht die Suppe und sagt: „Ist das Suppe oder \voc[das Gift]{Gift}{poison}?“ Fräulein Meier antwortet: „Lieber Historiker, halten Sie meine Suppe nicht für Gift.“ Der Briefträger lacht und sagt: „Ich halte sie für gut.“ Der Lehrer findet in der Tasche des Bauern ein kleines \voc[das Trinklied]{Trinklied}{drinking song}. Er liest es laut: „\voc[das Brüderlein]{Brüderlein}{little brother}, trink nicht allein; sitz bei uns und sei nicht klein.“ Fräulein Meier sagt: „Das ist kein gutes Lied, aber es ist freundlich.“ Der Briefträger singt das Trinklied \voc[noch einmal]{noch einmal}{once again}. Das Schäflein hört zu und schläft fast ein.

Am Nachmittag will der Bauer gehen. Das Schäflein steht aber bei Fräulein Meier und will nicht mitgehen. „Das Leben ist manchmal schwer“, sagt der Bauer. „Dieses Schäflein liebt die Küche mehr als den Stall.“ Fräulein Meier sagt: „Vielleicht bleibt es heute hier.“ Der Bauer fragt: „Ist das wirklich möglich?“ Der Briefträger antwortet: „Ja. Fräulein Meier hat ein großes Herz.“ Der Historiker sagt: „Und der Briefträger hat viel Liebe zu diesem Schäflein.“ Der Briefträger wird rot.

Gegen Abend hört der Regen auf. Der Himmel ist \voc{grau}{grey}, aber nicht finster. Fräulein Meier, der Briefträger, der Lehrer, der Bauer und der Historiker gehen in den Garten. Das Schäflein geht langsam hinter ihnen her. Auf der Erde liegt Wasser, und der Weg ist \voc{feucht}{wet, damp}. Am Ende des Gartens steht ein kleiner Stein. Von dort sieht man einen See. Der See \voc[spiegeln]{spiegelt}{to reflect, to mirror} den Himmel und die Wolken.

Am See ist ein kleiner \voc[der Meeresstrand]{Meeresstrand}{seashore}. Er ist nicht groß, aber Fräulein Meier nennt ihn gern ihren Meeresstrand. Auf dem Wasser liegt eine kleine \voc[die Insel]{Insel}{island}. „Die Insel ist \voc{einsam}{lonely}“, sagt Fräulein Meier. Der Briefträger antwortet: „Einsam ist sie nicht, denn eine Ente steht dort.“ Über dem Wasser beginnt die \voc[die Dämmerung]{Dämmerung}{twilight}. Alles wird still. Auch der Historiker \voc[schweigen]{schweigt}{to be silent} endlich.

Plötzlich hören sie einen \voc[der Ton]{Ton}{tone}. Der Ton kommt aus der \voc[die Tiefe]{Tiefe}{depth} des Gartens, vielleicht aus dem Busch. Dann hören sie eine \voc[die Stimme]{Stimme}{voice}. „Wer ruft dort?“ fragt der Briefträger. Der Lehrer sagt \voc{leise}{quietly}: „Vielleicht ist es das \voc[das Vogelrufen]{Vogelrufen}{bird call}.“ Ein Vogel \voc[fliegen]{fliegt}{to fly} über den See. Das Schäflein schaut auf den Vogel und bleibt stehen. Fräulein Meier sagt: „Das ist wie ein \voc[der Traum]{Traum}{dream}.“ Der Briefträger antwortet: „Ja, wie ein Traum ohne schlechte Nachricht.“

Auf dem Rückweg sagt der Historiker: „In einem alten Buch steht eine \voc[die Lüge]{Lüge}{lie}: Ein gutes \voc[das Weib]{Weib}{wife} spricht nie laut.“ Fräulein Meier bleibt stehen. „Das ist wirklich eine Lüge“, sagt sie. „Ein gutes Weib, eine gute Frau und auch ein gutes Fräulein dürfen laut sprechen.“ Der Briefträger lacht nicht, denn er will höflich sein. Der Lehrer sagt: „Dieser Satz ist besser als der alte Satz.“ Der Historiker nickt langsam.

Zu Hause ist es warm. Fräulein Meier sagt: „Bleiben Sie alle noch ein wenig.“ Der Bauer antwortet: „Ich muss gehen, aber das Schäflein darf heute bei Ihnen bleiben.“ Der Briefträger fragt: „Darf ich morgen noch einmal kommen?“ Fräulein Meier sieht ihn an und sagt: „Ja. Kommen Sie morgen noch einmal. Bringen Sie aber keinen Sack, keinen Bauern und keinen Bettler mit.“ Der Briefträger sagt: „Ich bringe vielleicht nur einen Brief.“ Fräulein Meier antwortet: „Das ist \voc[letzt]{letzten}{last} Endes genug.“

In der Nacht schläft das Schäflein in der Küche. Die Katze liegt auf dem Stuhl, der Hund schläft vor der Tür, und Fräulein Meier sitzt am Tisch. Sie schreibt in ihren Kalender: „Heute kam ein Schäflein. Es gehörte dem Bauern, aber einen Tag lang gehörte es auch zu meinem Haus.“ Dann denkt sie, dass das Leben manchmal arm und manchmal reich ist, dass Fragen nicht immer schwer sind, und dass ein guter Briefträger viel Freude bringt. Schließlich geht sie zu Bett. Das Haus ist still, und alle sind \voc[zu Hause]{zu Hause}{at home}.

\end{readersection}

\begin{readersection}{Das Abendbrot in der Schule}

Am Samstagabend saßen Jonas, Lara und Mira in der kleinen \voc[die Schule]{Schule}{school} am Fluß. Sie besuchten dort einen neuen \voc[der Kurs]{Kurs}{course} für deutsche Sprache. Draußen standen dunkle \voc[die Wolke]{Wolken}{cloud} über der Stadt, und die \voc[die Flut]{Flut}{flood} stieg schnell. Der Lehrer Herr Schneider sah aus dem Fenster und sagte: „Heute geht niemand allein nach Hause. Die Straße ist naß, und der Keller neben der Brücke ist vielleicht schon \voc{übervoll}{overflowing}.“

Auf dem langen \voc[der Tisch]{Tisch}{table} standen Brot, Käse, ein \voc[der Becher]{Becher}{cup} Wasser und ein kleines \voc[das Messer]{Messer}{knife}. Jonas dachte an das \voc[das Abendbrot]{Abendbrot}{supper} zu Hause. „\voc[mein]{Meine}{my} Mutter wartet sicher“, sagte er. Lara antwortete: „\voc[dein]{Deine}{your} Mutter weiß, \voc{daß}{that} du hier bist. Meine \voc[die Eltern]{Eltern}{parents} wissen es auch.“ Herr Schneider nickte: „Ja. Ihr bleibt hier, \voc{da}{since} der Weg am Fluß gefährlich ist.“

Lara war unruhig. Ihr kleiner Bruder Tom war noch nicht zurück. „Er spielt immer \voc{wie}{like} ein großer \voc[der Bote]{Bote}{messenger}“, sagte sie. „Er \voc[tragen]{trägt}{to wear, carry} einen Sack mit Briefen und will jedem helfen. Aber heute ist die Straße zu gefährlich.“ Jonas sagte: „Ich kenne ihn. Er ist vorsichtig.“ Lara wurde trotzdem \voc{böse}{angry, evil}: „Vorsichtig? Er ist zwölf!“ Herr Schneider antwortete ruhig: „Dein \voc[der Zorn]{Zorn}{wrath} hilft ihm nicht. Wir holen ihn, \voc{wenn}{if, when} er nicht bald kommt.“

In diesem Augenblick klopfte es. Ein alter \voc[der Gast]{Gast}{guest} trat ein. Sein Mantel war naß, und sein graues \voc[das Haupt]{Haupt}{head} zitterte. An der \voc[die Wand]{Wand}{wall} nahm Herr Schneider einen alten \voc[der Hirtenstab]{Hirtenstab}{shepherd's crook}. „Damit prüfe ich das Wasser“, sagte er. Der Gast zeigte zur Brücke: „Dort ist ein Mann mit einem Kind. Ich glaube, das Kind trägt einen Sack.“

„Tom!“ rief Lara. Herr Schneider sagte: „Jonas, du kommst mit mir. Lara, du auch. Mira bleibt hier.“ Mira sprang auf. „Mein Vater arbeitet in dem Haus neben der Brücke!“ „Gerade deshalb bleibst du“, sagte Herr Schneider. „Oskar bleibt bei dir.“ Oskar war ein stiller Junge. Paul nannte ihn manchmal einen \voc[der Horcher]{Horcher}{listener}, denn er hörte alles und sagte wenig.

Draußen gingen sie \voc{langsam}{slow, slowly} zur Brücke. Jonas trug die Lampe, Lara folgte ihm, und Herr Schneider ging voran. Paul wollte sie \voc[einholen]{einholen}{to catch up with}, aber Herr Schneider rief: „Bis hierher! Diese Mauer ist die \voc[die Grenze]{Grenze}{limit}. Niemand geht weiter ohne mich.“ Paul blieb stehen und \voc[folgen]{folgte}{to follow} ihnen nur mit den Augen.

Aus dem Haus kam Toms Stimme. „Hier!“ Unten im Keller stand Tom neben Miras Vater. Das Wasser war hoch. Tom hielt ein kleines \voc[das Tier]{Tier}{animal} im Arm, eine nasse Katze. „Sie saß auf einer Kiste“, sagte er. „Ich konnte sie nicht dort lassen.“ Lara nahm die Katze. „Gib sie mir.“ Die Katze biß fast in \voc[ihr]{ihren}{her, their} Ärmel, aber Lara hielt sie fest. „Sie hat Angst. Sie ist kein \voc[der Gegner]{Gegner}{opponent}.“

Miras Vater wollte noch eine Kiste retten. Herr Schneider sagte: „Lassen Sie die Kiste. Ihr Leben ist wichtiger.“ Der Mann sah auf seine Werkzeuge, dann auf Tom und Lara. „Gut“, sagte er. „\voc[außer]{Außer}{except for} dieser Tasche nehme ich nichts mit.“ Jonas nahm die Tasche. „Ich trage sie.“ Tom gab ihm den Sack mit Briefen. „Und den auch.“ Jonas lachte: „Heute \voc[kriegen]{kriege}{to get} ich viel Arbeit.“

Als sie zurückkamen, \voc[decken]{deckte}{to set, cover} Mira den Tisch. Oskar stellte die Becher hin. Der alte Gast fragte \voc[fragen nach]{nach}{to ask about} Tom und sagte dann zu Herr Schneider: „\voc[Ihr]{Ihr}{your} Kurs ist heute anders als gewöhnlich.“ Herr Schneider lächelte: „Ja. Heute lernen wir nicht nur Wörter. Heute lernen wir, daß eine gute \voc[die Tat]{Tat}{deed} auch Verstand braucht.“

Tom wurde rot. „Ich wollte nur helfen.“ Lara setzte sich neben ihn. „Das war gut. Aber du sollst nicht allein zur Flut gehen.“ Tom zog ein nasses Papier aus der Tasche. „Dann gebe ich euch ein \voc[das Versprechen]{Versprechen}{promise}: Bei Flut gehe ich nicht allein.“ „\voc{Wahrlich}{verily}“, sagte der alte Gast, „das ist ein gutes Versprechen.“

Mira saß bei ihrem Vater. Sie hatte viel \voc[der Kummer]{Kummer}{trouble} gehabt, aber jetzt fühlte sie \voc[die Freude]{Freude}{joy}. „Du willst immer alles retten“, sagte sie zu ihm. „Das ist keine \voc[die Faulheit]{Faulheit}{laziness}, aber manchmal auch keine Weisheit.“ Ihr Vater nickte. „Ich wollte meine Sachen \voc[sparen]{sparen}{to save}. Doch heute habt ihr mich gerettet.“

Die Katze lag unter dem Tisch. Neben ihr kroch eine kleine \voc[die Ameise]{Ameise}{ant} über den Boden. Tom wollte sie wegschieben, aber Lara sagte: „Laß sie. Heute retten wir auch kleine Dinge.“ Herr Schneider sah die Schüler an. „Jeder hat eine \voc[die Fähigkeit]{Fähigkeit}{ability}: einer trägt, einer hört, einer wartet, einer denkt. Das ist auch eine \voc{geistig}{spiritual, intellectual}e Arbeit.“

Draußen fiel das Wasser langsam. „Die \voc[die Ebbe]{Ebbe}{ebb} kommt“, sagte Jonas. „Dann können wir bald nach Hause.“ Lara sah zu Tom. „Aber du gehst \voc[nicht]{nicht}{not} allein.“ Tom antwortete: „Nein. Ich gehe mit dir.“ Jonas hob seinen Becher. „Auf \voc[unser]{unseren}{our} Kurs.“ Der alte Gast lachte: „Und auf die Katze.“ Tom sah unter den Tisch. „Sie hat ihren \voc[eigen]{eigenen}{own} Kopf.“ Lara sagte: „\voc{Doch}{nevertheless} bleibt sie vielleicht bei uns.“ Alle lachten, und aus dem Abendbrot in der Schule wurde eine Geschichte, die sie lange nicht vergaßen.

\end{readersection}

\begin{readersection}{Der Ruhetag am Fluß}

\voc[heute]{Heute}{today} ist in der kleinen Stadt kein \voc[der Werktag]{Werktag}{work day}, sondern ein \voc[der Feiertag]{Feiertag}{holiday}. Für viele Leute ist er auch ein \voc[der Ruhetag]{Ruhetag}{day of rest}. Jonas und Lara gehen früh an den Fluß. Über dem Wasser liegt grauer \voc[der Rauch]{Rauch}{smoke} aus den Häusern, und am Ufer sehen sie eine tiefe \voc[die Fußspur]{Fußspur}{footprint} im feuchten Boden.

„Das ist keine Spur von einem Hund“, sagt Lara. „Vielleicht gehört sie einem großen Tier.“ Jonas kniet nieder. Er will \voc{erkennen}{to recognize}, welches Tier hier gewesen ist. „Hier sieht man keine \voc[die Klaue]{Klauen}{claw}“, sagt er. „Also war es wohl kein Wolf und kein \voc[der Löwe]{Löwe}{lion}.“ Lara lacht. „Ein Löwe in unserer Stadt? Das wäre eine gute Nachricht für die Zeitung.“

Auf der Brücke steht ein alter \voc[der Grieche]{Grieche}{Greek} mit einem Buch unter dem Arm. Jonas kennt den Griechen aus der Schule, denn der Grieche zeigt den Schülern manchmal alte Bücher. Das Buch sieht \voc{antik}{ancient} aus; ein dünner \voc[der Faden]{Faden}{thread} hält die Blätter zusammen. Auf der ersten Seite steht ein \voc[aktuell]{aktuelles}{current, modern} Datum, aber darunter stehen viel ältere Worte.

Lara \voc[fragen nach]{fragt nach}{to ask about} dem Buch. Der Grieche öffnet es vorsichtig. „Es \voc[gehören]{gehört}{to belong to} meiner Familie“, sagt er. „Mein Großvater schrieb hier seine Notizen hinein. Er war ein \voc[fleißig]{fleißiger}{industrious} Lehrer, aber kein reicher Mann. Er schrieb über Tiere, Wasser, Pflanzen und den \voc[der Mensch]{Menschen}{person}.“ Jonas liest eine Zeile und sagt: „Hier steht: Der Mensch erkennt die Welt nicht an einem Tag.“ Der Grieche nickt. „Das war sein liebster \voc[der Gedanke]{Gedanke}{thought}.“

Am Ufer sitzt ein junger Mann auf einem Stein. Er sieht \voc{krank}{sick} aus, und an seiner Hand ist eine kleine \voc[die Wunde]{Wunde}{wound}. Lara geht zu ihm. „Brauchen Sie Hilfe?“ Der Kranke hebt den \voc[der Kopf]{Kopf}{head}. „Nein, es ist \voc{nichts}{nothing} Schlimmes. Ich wollte nur an den Fluß gehen. Der Weg war länger, als ich dachte.“ Jonas gibt ihm Wasser. „Es ist nicht warm, aber es ist \voc{rein}{pure}.“ Der Kranke trinkt und lächelt. „Danke. Das ist \voc{angenehm}{agreeable}.“

Der Grieche setzt sich zu ihnen. „Ein Kranker soll an einem Feiertag nicht allein gehen“, sagt er. Lara antwortet: „Ein \voc[gesund]{Gesunder}{healthy} auch nicht immer.“ Der Kranke lacht leise. „Dann bin ich heute gut begleitet.“ Jonas sieht die alte Spur im Boden und sagt: „Vielleicht war das Tier auch krank.“ Lara sieht ihn an. „Oder es war nur müde.“ Der Grieche sagt: „\voc{Manchmal}{sometimes} ist ein kleiner Weg für einen müden Körper schon sehr lang.“

Plötzlich läuft eine Katze aus dem Busch. Sie trägt einen kleinen Stein im Maul und läßt ihn vor Laras Fuß fallen. Der Stein ist glatt und fast weiß. „Sie bringt uns ein Geschenk“, sagt Jonas. Lara hebt den Stein auf. „Das Reine ist manchmal klein“, sagt sie, aber dann lacht sie gleich. „Entschuldigung, das klingt zu sehr wie ein Sprichwort.“ Der Kranke sagt: „Nein, heute paßt es.“

Ein Junge kommt den Weg entlang. Er sieht \voc{elend}{miserable} aus, denn seine Schuhe sind naß und sein Brot ist ins Wasser gefallen. „Ich wollte es der Katze geben“, sagt er. „Jetzt hat sie nichts.“ Lara gibt ihm ein Stück Brot. „Dann teilen wir.“ Der Junge will das Brot sofort zur Katze \voc{werfen}{to throw}, aber Jonas hält seine Hand fest. „Nicht werfen. Gib es ihr langsam.“ Der Junge tut es, und die Katze kommt vorsichtig näher.

„Du bist \voc{kühn}{daring, bold}“, sagt der Grieche zu dem Jungen. „Aber nicht jede kühne Tat ist auch klug.“ Der Junge runzelt die Stirn. „Was heißt das?“ Der Grieche zeigt auf die Katze. „Wenn du das Brot wirfst, läuft sie vielleicht weg. Wenn du wartest, kommt sie zu dir.“ Der Junge denkt nach. „Dann ist Warten auch eine Tat?“ Der Grieche lächelt. „Vielleicht. Ein \voc[weise]{weiser}{wise, sage} Mensch weiß das.“

Lara sieht den Griechen an. „Sind Sie also ein Weiser?“ Der Grieche lacht. „Nein. Ich bin nur ein alter Grieche mit einem alten Buch.“ Der Kranke sagt: „Aber dem Weisen hört man gern zu.“ Jonas merkt sich den Satz, denn in der Schule hatten sie gerade gelernt, daß man aus Adjektiven Namen machen kann: der Kranke, der Gesunde, der Weise. Das ist schwer, aber es hilft ihm beim Lesen.

Am Mittag kommen andere Leute an den Fluß. Sie sind \voc{verschieden}{different}: ein Arbeiter, zwei Kinder, eine alte Frau und ein Student. Der Student will \voc{alles}{everything} wissen: Wem gehört das Buch? Wer hat die Spur gesehen? Ist der Kranke verletzt? Woher kommt der Grieche? Lara sagt: „Du fragst sehr viel.“ Der Student antwortet: „Ich schreibe einen kleinen Bericht. Vielleicht kann er in die Stadtzeitung \voc{hineinpassen}{fit in}.“

Da kommt der Arbeiter wieder. Neben ihm läuft ein junges Kalb. „Das war kein Löwe“, sagt er. „Das war mein Kalb. Es ist heute früh weggelaufen.“ Alle lachen. Jonas zeigt auf die Spur. „Dann haben wir viel gedacht und wenig gefunden.“ Der Grieche antwortet: „Dennoch war der Gedanke nicht verloren.“ Der Arbeiter hält das Kalb fest. „Der \voc[der Wille]{Wille}{will} dieses Tieres ist stark“, sagt er. „Aber mein Wille ist heute stärker.“

Am Nachmittag wird der Himmel dunkel, und der Wind weht \voc{gewaltig}{powerfully}. Der Rauch über dem Wasser verschwindet, und der Fluß sieht still aus. Der Kranke kann wieder gehen. Lara begleitet ihn bis zur Brücke. „\voc[morgen]{Morgen}{tomorrow} ist wieder ein Werktag“, sagt er. „Dann muß ich vorsichtiger sein.“ Lara antwortet: „Und Sie dürfen den Kranken in sich nicht \voc{verachten}{to despise}.“ Der Mann sieht sie überrascht an. „Das klingt wieder wie ein Sprichwort.“ Lara lacht. „Vielleicht. Aber diesmal meine ich es einfach.“

Der Grieche gibt Jonas das alte Buch. „Trag es bis zur Schule“, sagt er. „Der Faden ist schwach.“ Jonas nimmt es vorsichtig. „Ich möchte es nicht verlieren.“ Der Grieche nickt. „Ein Buch ist manchmal wie ein Mensch. Es braucht Hilfe, auch wenn es alt ist.“ Dann zeigt er auf eine kleine Zeichnung in dem Buch. Sie zeigt einen \voc[heilig]{Heiligen}{holy, saintly} mit einem Tier. „Für meinen Großvater war nicht jedes alte Wort heilig“, sagt er, „aber er suchte das Heilige im einfachen Leben.“

\voc[abends]{Abends}{evenings} sitzt Jonas zu Hause am Tisch. Sein Kopf ist müde, aber nicht leer. Er schreibt in sein Heft: „Ein Ruhetag ist nicht immer ruhig.“ Dann schreibt er darunter: „Der Kranke, der Gesunde, der Weise, der Grieche und der Mensch: alle tragen ihre eigene Geschichte.“ Der \voc[der Zeiger]{Zeiger}{indicator} der Uhr steht auf zehn. Jonas legt den Stift weg. Draußen ist die Stadt still, und in seinem Heft bleibt ein kleiner \voc[der Wunsch]{Wunsch}{wish} zurück: Am nächsten Feiertag will er wieder an den Fluß gehen.

\end{readersection}

\begin{readersection}{Der Dieb in der Küche}

In einer kleinen Stadt gab es ein altes Gasthaus. Es stand nicht \voc{weit}{far} von der Brücke, neben einem großen \voc[der Baum]{Baum}{tree}. Am Morgen war der \voc[der Himmel]{Himmel}{heaven, sky} oft grau, aber in der \voc[die Küche]{Küche}{kitchen} war es warm. Dort arbeitete ein junger \voc[der Koch]{Koch}{cook}. Er hieß Lukas, und die Leute der Stadt hatten ihn gern, denn er kochte gut und lachte viel.

Seit drei Tagen aber hatte Lukas \voc[das Kopfweh]{Kopfweh}{headache}. In der Küche war etwas nicht in Ordnung. Am Abend lagen fünf Bratwürste auf dem Tisch; am Morgen waren es \voc{nur}{only} vier. Ein Kuchen stand neben dem Fenster; später fehlte ein großes Stück. Eine alte Frau sagte: „\voc[es gibt]{Es gibt}{there is, there are} einen \voc[der Dieb]{Dieb}{thief} in der Küche.“ Ein Kind sagte: „Vielleicht ist es ein kleiner böser Mann.“ Ein anderes Kind sagte: „Vielleicht ist es ein \voc[ander]{anderer}{other} Mann, der sehr hungrig ist.“ Ein \voc[gelehrt]{gelehrter}{learned, scholarly} Herr aus der Schule meinte: „Vielleicht ist es kein Mensch, sondern ein Tier.“

Lukas schrieb deshalb eine \voc[die Regel]{Regel}{rule} an die Tür: „Niemand \voc[gehen]{geht}{to go} nachts in die Küche.“ Darunter schrieb er eine zweite Regel: „\voc{Auch}{also, even} kein Hund.“ Die dritte Regel war noch strenger: „Auch keine Katze.“ Aber am nächsten Morgen fehlte wieder ein Stück Apfelkuchen. Lukas sah den leeren Teller und sagte traurig: „Diese Regeln helfen nicht. Es gibt immer eine \voc[die Ausnahme]{Ausnahme}{exception}.“

In der Stadt wohnte ein alter Richter. Er hörte von dem Dieb und kam am Mittag ins Gasthaus. Er trug ein schwarzes Kleid und sah sehr ernst aus. „Ein Dieb muß ins \voc[das Gefängnis]{Gefängnis}{prison}“, sagte er. „Aber zuerst muß man ihn finden.“ Lukas zeigte ihm die Küche, den Tisch, das Fenster und die Tür. Der Richter sah unter den Tisch, hinter die Tür und in den kleinen Sack mit Mehl. Dort war niemand.

„Vielleicht kommt der Dieb durch das Fenster“, sagte der Richter. „Es ist niedrig genug.“ Lukas sah hinaus. Vor dem Fenster stand der Baum. Ein langer Ast ging fast bis zum Dach. „Da kann ein Mensch nicht gehen“, sagte Lukas. „Aber ein Tier kann vielleicht von dem Baum auf das Dach und dann zum Fenster kommen.“

Am Abend blieb Lukas allein in der Küche. Er stellte Brot, Käse und eine Bratwurst auf den Tisch. Dann löschte er das Feuer und setzte sich hinter einen großen Sack. Es wurde dunkel. Draußen \voc[fallen]{fiel}{to fall} Regen, und es war still. Eine Uhr schlug neun, dann zehn. Lukas hatte wieder Kopfweh, aber er wartete.

Endlich hörte er ein Geräusch. Etwas fiel leise vom Baum auf das Dach. Dann kam ein kleines Gesicht am Fenster hervor. Es war kein Mensch. Es war kein Hund. Es war eine magere, graue Katze mit einem \voc[häßlich]{häßlichen}{ugly} Ohr. Sie sprang in die Küche, ging langsam zum Tisch und \voc[beißen]{biß}{to bite} in die Bratwurst. Lukas rief: „Aha! Jetzt sehe ich den Dieb!“ Die Katze erschrak, wollte weglaufen, aber dabei fiel ein altes \voc[das Hufeisen]{Hufeisen}{horseshoe} aus dem Fenster auf den Boden.

Am nächsten Morgen kamen viele Leute ins Gasthaus. Der Richter, der gelehrte Herr, zwei Kinder und eine Frau mit roter \voc[die Krone]{Krone}{crown} auf dem Hut standen in der Küche. „Das ist keine Krone“, sagte ein Kind. „Das ist nur ein Hut.“ Die Frau lachte. Der Richter aber sah sehr ernst auf die Katze. „Du hast gestohlen“, sagte er. „Aber du bist kein Mensch. Was macht man mit einer solchen Diebin?“

Die Katze saß unter dem Tisch und sah Lukas an. Sie hatte Angst, und sie war sehr dünn. „\voc[solang(e)]{Solang}{as long as} sie Hunger hat, kommt sie wieder“, sagte Lukas. Der gelehrte Herr nickte. „Das ist richtig. Solang der Bauch leer ist, fragt das Tier nicht nach Regeln.“

Da sagte die alte Frau: „Man soll die Katze nicht töten und nicht verjagen. Es wäre \voc[tödlich]{tödlich}{deathly} für sie, wenn sie im Winter draußen bleibt.“ Der Richter antwortete: „Niemand will, daß sie draußen \voc{sterben}{to die} muß. Aber die Ehre des Gasthauses muß bleiben.“ Lukas dachte nach. „Die \voc[die Ehre]{Ehre}{honor} des Gasthauses wird nicht kleiner, wenn eine arme Katze etwas Essen bekommt.“ Die alte Frau sagte leise: „Auch eine Katze hat ihre kleine \voc[die Würde]{Würde}{dignity}.“

Die Kinder lachten. „Dann ist sie also kein Dieb mehr?“ Der Richter sagte: „Vielleicht \voc[wechseln]{wechselt}{to change} sie ihren Beruf.“ Lukas stellte eine kleine Schale Milch unter den Tisch. „Milch ist vielleicht kein \voc[das Gegengift]{Gegengift}{antidote} gegen Hunger“, sagte er, „aber sie hilft besser als ein Gefängnis.“ Die Katze kam vorsichtig hervor und trank. Dann legte sie sich neben den warmen Ofen. „Jetzt gibt es eine neue Regel“, sagte Lukas. „Die Katze darf hier bleiben, aber sie darf nicht auf den Tisch springen.“

Am Nachmittag kam der Herrscher der Stadt in das Gasthaus. Eigentlich war er kein wirklicher König, sondern nur der Bürgermeister; aber die Kinder nannten ihn den \voc[der Herrscher]{Herrscher}{ruler}, weil er bei Festen eine goldene Kette trug. Er hörte die Geschichte und sagte: „Ein Gasthaus mit einer Katze ist besser als ein Gasthaus mit einem Gefängnis.“ Dann hob er das alte Hufeisen auf. „Das ist ein Zeichen des Glückes. Hängt es über die Tür.“

Lukas hängte das Hufeisen über die Küchentür. Es sah ein wenig häßlich aus, aber die Leute fanden es lustig. Von diesem Tag an gab es im Gasthaus keine gestohlenen Bratwürste mehr. Es gab nur eine graue Katze, die jeden Abend unter dem Tisch schlief. Die Kinder nannten sie „mein \voc[das Liebchen]{Liebchen}{little loved one}“, obwohl sie manchmal noch immer biß. Wenn jemand ihr \voc[zu]{zu}{too} nahe kam, zeigte sie die Zähne.

Der Richter kam später oft zum Essen. Er sagte jedes Mal: „Die Regel ist wichtig.“ Lukas antwortete: „Ja, aber die Ausnahme ist manchmal lebendig.“ Dann lachten die Gäste. Draußen stand der alte Baum, der Himmel wurde dunkel, und in der warmen Küche schlief die Katze unter dem Tisch, nicht weit vom Feuer.

\end{readersection}


\end{document}
